\documentclass[10pt,a4paper]{amsart}
\usepackage[margin=19mm]{geometry}
\usepackage{amsmath,amssymb,mathtools}
\usepackage[scr=rsfs,cal=euler]{mathalfa}
\usepackage{xfrac}
\usepackage[unicode,hidelinks,bookmarks=false]{hyperref}
\newcommand{\ie}{i.e.\ }
\newcommand{\cf}{cf.\ }
\newcommand{\resp}{respectively\ }
\newcommand{\Subsection}{\subsection}
\providecommand{\nobreakdash}{\nobreak\hbox{-}}
\setlength{\emergencystretch}{2em}
\multlinegap0pt
\usepackage{amssymb}
\newtheorem{lemma}{Lemma}
\title{Every Finite Group Admits a Just Finite Presentation}
\author{Marc Lackenby}
\date{Source: arXiv:2605.10402v1 (CC BY 4.0)}
\begin{document}
\maketitle

\noindent\emph{Source and licence.} Every Finite Group Admits a Just Finite Presentation. Version arXiv:2605.10402v1 (CC BY 4.0), distributed by arXiv under the Creative Commons Attribution 4.0 International licence, http://creativecommons.org/licenses/by/4.0/, as recorded in the arXiv licence metadata for this exact version and shown on the abstract page https://arxiv.org/abs/2605.10402v1. Full licence notice: \texttt{licenses/arxiv-2605.10402-CC-BY-4.0.txt}.\par\smallskip

\noindent\textit{Editorial scope.} Counted unit: the article's lemma Lem:SemiDirect (source0.tex lines 73--76). The article's complete original proof of it is delivered with it (source0.tex lines 78--80, one \texttt{proof} environment of 3 lines). No mathematics is written, repaired or re-derived: the statement and the proof are the article's own lines, in the article's own notation, and no retained line is substituted or re-flowed.  No further source-local context is needed: every cross-reference of every retained line resolves inside the two delivered spans.  Nothing was omitted from any retained span, so the package carries no omission record.  No retained line contains a citation, so the wrapper carries no bibliography and no bibliography entry.  No retained line contains a figure environment, an image-inclusion command or drawing code, so no image is delivered and the package declares no asset dependency.  Editorial formatting only, in addition: an AMS \texttt{amsart} preamble that restates the article's own declarations and macros used by the retained lines, namely the environments \texttt{lemma} on separate counters, together with the standard packages those lines need.  The retained environments print in this document as Lemma 1 in order of appearance, whereas the article prints the same items with the numbering of its own full document.  The complete native source of the article as distributed in the arXiv e-print for this exact version is bundled unchanged as \texttt{sources/200294/source0.tex}.
\begin{lemma}
\label{Lem:SemiDirect}
For $k > 1$, the presentation $\langle x, y \, | \, x^{-1} y x = y^2, x^k = 1 \rangle$ specifies the group $\mathbb{Z}_{2^k-1} \rtimes \mathbb{Z}_k$, where $\langle x \rangle = \mathbb{Z}_k$ and $\langle y \rangle = \mathbb{Z}_{2^k-1}$ and the conjugation action of $x$ on $\langle y \rangle$ is $y \mapsto y^2$.
\end{lemma}

\begin{proof}
There is a homomorphism $\langle x, y \, | \, x^{-1} y x = y^2, x^k = 1 \rangle \rightarrow \mathbb{Z}_{2^k-1} \rtimes \mathbb{Z}_k$, since the two relations are satisfied in the image group. It is surjective because the images of $x$ and $y$ generate $\mathbb{Z}_{2^k-1} \rtimes \mathbb{Z}_k$. Any element in $\langle x, y \, | \, x^{-1} y x = y^2, x^k = 1 \rangle$ can be written as $x^m y^n$ for some integers $m$ and $n$. Using $x^k = 1$, we may assume $0 \leq m < k$. Conjugating $y$ by $x$ a total of $k$ times gives $y^{2^k}$. Since $x^k = 1$, we deduce that $y^{2^k} = y$. So, we may also assume $0 \leq n < 2^{k}-1$. So $\langle x, y \, | \, x^{-1} y x = y^2, x^k = 1 \rangle$ has at most $k(2^k -1)$ elements. However, $\mathbb{Z}_{2^k-1} \rtimes \mathbb{Z}_k$ has exactly this many elements, and therefore the above surjective homomorphism is an isomorphism.
\end{proof}

\end{document}
